\chapter{华为云环境信息}

\section{设计目标与总体架构}

本课程设计基于华为云搭建一套完整的云原生平台：以云容器引擎 CCE 托管的
Kubernetes 集群为核心，部署「Flask 后端 API + Redis 数据库」两层 Web 应用，
依次实现容器化、对外暴露、配置分离、数据持久化与弹性伸缩；随后在同一集群上
通过 Spark Operator 运行 PySpark 大数据分析作业；最后扩展完成监控系统、
CI/CD 流水线与分布式 AI 训练三个附加专题。总体技术架构如图~\ref{fig:arch} 所示。

\begin{figure}[htbp]\centering
\begin{minipage}[c][5.2cm][c]{0.9\textwidth}\centering\ttfamily\small
\fbox{\parbox{0.95\linewidth}{\centering\vspace{2pt}
开发者 / GitHub Actions \\[2pt]
$\Downarrow$ docker build / push \\[2pt]
SWR 私有镜像仓库 (group17) \\[2pt]
$\Downarrow$ image pull \\[4pt]
\rule{0.9\linewidth}{0.4pt}\\[2pt]
CCE Kubernetes 集群（2 Worker 节点）\\[3pt]
ELB(公网) $\to$ backend(Deploy×2) $\to$ redis(PVC) \\[2pt]
frontend(Nginx) ｜ HPA ｜ Spark Operator ｜ Prometheus/Grafana \\[3pt]
\rule{0.9\linewidth}{0.4pt}\\[2pt]
OBS 对象存储 (group17)：豆瓣电影数据集 (s3a://)
\vspace{2pt}}}
\end{minipage}
\caption{基于华为云的云原生平台总体架构}\label{fig:arch}
\end{figure}

\section{华为云服务与账号信息}

本设计使用了表~\ref{tab:hwservices} 所列的华为云服务，区域（Region）统一为
\textbf{华北-北京四（cn-north-4）}。

\begin{table}[htbp]\centering
\caption{华为云服务与命名}\label{tab:hwservices}
\begin{tabular}{lll}
\toprule
服务 & 作用 & 本设计取值 \\
\midrule
CCE  & 托管 Kubernetes 集群 & 集群 group17，版本 v1.35.3 \\
SWR  & 私有镜像仓库 & 组织 group17 \\
OBS  & 对象存储（数据集） & 桶 group17（s3a 接口） \\
ELB  & 弹性负载均衡（公网入口） & 自动创建，公网 IP 1.94.238.112 \\
\bottomrule
\end{tabular}
\end{table}

\section{CCE 集群与节点规格}

集群由华为云托管 Master 节点（免运维），学生创建并管理 2 个 Worker 节点，
节点规格为 \textbf{c9.large.4（2~vCPU / 8~GiB）}，操作系统为 Huawei Cloud
EulerOS 2.0，容器运行时为 containerd~2.2.1，Kubernetes 版本 v1.35.3
（$\geq 1.27$，满足要求）。节点信息如图~\ref{fig:nodes}（任务2 验收截图）所示。

\shot{task2_get_nodes.png}{kubectl get nodes -o wide：2 个 Worker 节点 Ready，版本 v1.35.3}{fig:nodes}

为支撑后续 HPA（任务6）与监控，集群额外部署了 metrics-server（提供
\texttt{metrics.k8s.io} 资源指标 API），\texttt{kubectl top nodes} 可正常读取
节点 CPU/内存用量。
